\section{Transferring the three-command compiler to the unary cover calculus}
\label{sec:transfer-u}

This argument concerns the changed register source, independently of physical
allocation.  Let $\mathcal D_U$ be the exact numerical calculus calibrated by
the primary U214 report: truth $0$, falsity $1$, the seven retained logical
rows, dominated propositional row, modus ponens, theorem generalization,
and the enabled role-reassigned cover bundle $B^*(a,b,c)$.  The role-emitter
calibration of U214 supplies, as a mathematical dependency, the
PA-formalizable equivalence
\[
 \mathcal D_U\vdash 1\quad\longleftrightarrow\quad
 \neg\operatorname{Con}(U)\quad\longleftrightarrow\quad
 \neg\operatorname{Con}(\mathsf{PA}).
\]
The K215 transfer compiler can be used here because its command proof depends
on the generated rows' theoremhood, their numerical tags, their selectable
positions, and the gate, rather than on the set-theoretic content of the
old covering bundle.  The new bundle satisfies all of these interfaces.
No theorem about the literal ZFC machine is identified with a theorem about
the changed arithmetic machine.

Write $P(x,y)=(x+y)(x+y+1)/2+y$, $I(x,y)=4P(x,y)$,
$Q_i(x)=4P(i,x)+2$, and $[x;w]=P(x,w)$.  For a command store
$(t,w,d,h)$, define the direct transfer
\[
 S(t,[x;v],d,h)=
 \begin{cases}\mathsf{halt},&h=1,\\(h,v,d,x),&h\ne1.\end{cases}
\]
There is no incoming-top guard in this definition.  The unchanged AX
arithmetic has, with $A=d+x$, $B=h+y$, $C=c$, and $U=d$, the effect
\[
 \operatorname{AX}(c;x,y)(t,w,d,h)=
 (\operatorname{gate}(A,B,c),
 [r_1,\ldots,r_7,\ell(A,B,c,U),Q_U(h),t;w],0,0).
\]
The gate is $B^*(A,B,c)$ when $c\ne A,B$ and is zero otherwise.
The eight logical rows and the gate are implication codes and hence are
not $1$.  The quantified entry has residue $2$ modulo $4$, so it is
neither $0$ nor $1$.  These assertions hold for every natural choice of
parameters, including the alias $I(0,0)=0$.

\subsection{The two modus-ponens implementations}

Both versions compute $Y=d+x$ and compare $I(t,Y)$ with $h$.  Failure
returns $(0,w,Y,0)$.  Ordinary MP returns $(Y,[0;w],Y,0)$ on success;
cache-return MP instead returns $(0,w,0,Y)$.  For the latter, the exact
source replacement is the closed loop
\[
 (\mathrm{CD}_4,(\mathrm{CI}_5,\mathrm{CC}_2)),
\]
which moves the target into the known-zero cache and clears the additive
counter.  This loop is justified for arbitrary counter values by induction.
The direct $S$ implementation is clear-top; copy-cache-to-top; old-NIP.
Its components have $2$, $7$, and $26$ primitive occurrences.  The copy
restores the cache and zero scratch, allowing old-NIP to test that cache
and then pop the numerical stack.  Its two associations have the same
all-natural contract.

\subsection{Soundness}

At a safe boundary, $t$, $h$, and every explicit stack head are
$\mathcal D_U$-theorems.  AX preserves this property: its ordinary rows
are axioms, its dominated row meets $U\le A$, its quantified head is
generalization of a theorem, and its gate is an axiom or truth.  Failed
MP returns truth.  Successful MP applies numerical modus ponens and stores
the conclusion either as current top or as cache.  Transfer moves theorem
codes and halts only when the cache is $1$.  Thus every halt from a safe
seed yields a numerical refutation.  This proves soundness for arbitrary
decoded fields, without a well-formedness or proof-word filter.

\subsection{Strong completeness}

For every theorem code $X$, the required compiler invariant is
\[
 (t,w,0,0)\xrightarrow{B_X}\mathsf{halt}
 \quad\text{or}\quad (X,w,0,0)
 \qquad\text{for every natural }t,w.
\]
This is deliberately stronger than the safe-execution invariant; the
retained top and stack need not be theorem-valued in this completeness
statement.

Put $D=\operatorname{AX}(0;0,0)$ and define
\[
 H=D;S^{11},\qquad
 \mathrm{PUSH}=D;S^9;\mathrm{MP}(0),\qquad
 \mathrm{DISCARD}=\mathrm{MP}(1);\operatorname{AX}(1;0,0);S^{10};\mathrm{MP}(0).
\]
The exact contracts are as follows.  $H$ tests the incoming top, and
otherwise preserves it, pops one preexisting stack head, and clears $d$.
PUSH never halts and returns $(0,[t;w],0,0)$.  DISCARD, from a zero
counter and cache, never halts and returns $(0,w,0,0)$.
The MP uses inside PUSH and DISCARD are forced failures.  In PUSH its
cache has residue $2$, whereas the compared implication has residue $0$.
In DISCARD the first comparison has $I(t,1)>0$ against cache zero;
the last has $I(Q_1(0),0)>0$ against zero.  Consequently these contracts
hold for both MP implementations.  The first ten transfers of $H$ test
only zero and generated nonfalsity entries; its eleventh tests the
original top.  The original cache is overwritten by AX before any test.

From top, counter, and cache zero, set $J(0)$ to the empty word and
$J(u)=\mathrm{MP}(u)$ for $u>0$.  This sets the counter to $u$ without
altering the stack.  The zero exception is essential: MP$(0)$ would
succeed by the raw zero alias.

Truth is compiled by DISCARD.  For $1\le j\le7$, a row $r_j(a,b,c)$
is compiled by
\[
 \mathrm{DISCARD};\operatorname{AX}(c;a,b);S^{j+1};H^{9-j}.
\]
For the dominated row with $u\le a$ use
\[
 \mathrm{DISCARD};J(u);\operatorname{AX}(c;a-u,b);S^9;H.
\]
For an enabled bundle use
\[
 \mathrm{DISCARD};\operatorname{AX}(c;a,b);H^{10}.
\]
All discarded entries, including the saved zero, are removed, while the
selected theorem is retained and the original stack is restored.
For generalization use
\[
 B_{Q_i(X)}=B_X;\mathrm{PUSH};J(i);S;\operatorname{AX}(i;0,0);S^{10}.
\]
Just before AX the tuple is $(0,w,i,X)$; the gate is disabled because
$A=c=i$, and ten transfers select $Q_i(X)$ and restore the stack.
This calculation includes $i=0$ and $X=1$.

For ordinary MP, let the premises be $A$ and $M=I(A,Y)$.  Use
\[
 B_M;\mathrm{PUSH};B_A;\mathrm{PUSH};S;S;\mathrm{MP}(Y);
 \operatorname{AX}(Y;0,0);S^{11}.
\]
Absent an earlier successful halt, the two transfers leave
$(A,w,0,M)$, MP yields $(Y,[0;w],Y,0)$, and the last AX/transfer
block clears $d$ and removes precisely that pushed zero.  Its final
transfer halts if $Y=1$; otherwise it returns $(Y,w,0,0)$.

For cache-return MP, reserve a disposable zero before computing either
premise:
\[
 \mathrm{DISCARD};\mathrm{PUSH};B_M;\mathrm{PUSH};B_A;\mathrm{PUSH};
 S;S;\mathrm{MP}(Y);S.
\]
The two shifts now leave $(A,[0;w],0,M)$, successful MP yields
$(0,[0;w],0,Y)$, and the last transfer either halts at $Y=1$ or
returns $(Y,w,0,0)$.  Reserving zero is necessary to avoid popping
an arbitrary head of the original stack.  The pairing identities make
all these traces valid when $A=Y=0$ or $w=0$, notwithstanding
$P(0,0)=0$.  Induction on a finite numerical derivation proves the
strong contract in both cases; $B_1;H$ then halts from every
$(t,w,0,0)$.

The ordinary selector fields are $S=[0]$, MP$(x)=[1,x]$, and
AX$(c;x,y)=[c+2,y,x]$.  Exchanging the two small command bodies
changes only their first selectors to $1$ and $0$, respectively.
Let $R$ be the successor-delimited continuation code of $B_1;H$.
For every $z$ the candidate $P(2,P(z,P(1,R)))$ first executes
AX$(0;1,z)$, clearing counter and cache, then executes the successful
suffix.  Since this candidate is at least $z$, successful candidates
are cofinal.  Every nonhalting continuation decreases after its
predecessor operation and Cantor projections, so the ordinary scheduler
reaches such a candidate beyond every finite initial cutoff.

\subsection{Startup dependency of the new table}

The actual shortened initializer is calibrated by the interior-AX
suffix lemma in Section~\ref{sec:u213-bootstrap}.  It reaches source
ordinal $136$ after $48{,}140$ binary steps, matching $141$ source
steps from ordinal $106$ and store $(1,0,4,0,0,1,0,0,1)$.
That interior suffix yields nine explicit theorem heads above opaque
tail $4$, next candidate $1$, and empty proof pointer.  The exact
three-command credit is $d(1)=2$, leaving permanent margin seven.
The old ST215 whole-AX startup is historical background and is not
used as a premise for this table.  The literal frame and whole-tape
primitive simulation are separate machine-specific checks.

The source/compiler argument uses finite syntax, natural-number
induction, and primitive-recursive proof transformations.  Its
PA-formalizable conclusion uses the inherited U214 arithmetic
calibration; no newly kernel-checked PA endpoint is asserted here.
